% 第 2 章　单位是物理世界的语法
\chapter{单位是物理世界的语法}

上一章说过，物理学从测量开始。可是你有没有想过：测量得到的那串数字，单独拿出来其实什么也不是。厨房秤上显示“500”，是 500 克，还是 500 毫克，还是 500 斤？数字后面那个小尾巴——单位——才是这句话真正的谓语。这一章我们暂时不谈任何新定律，只谈一件看似琐碎、实则致命的小事：物理世界里的每个数字都必须带着单位说话，而单位背后藏着一整套语法。学会这套语法，你就能在不解方程的情况下看出一个公式的错误，在只有常识的情况下估出一座城市的用水量。

\step{反常识开场}

1999 年 9 月 23 日，美国宇航局的火星气候轨道器在飞行了九个多月、跨越六亿多公里之后，终于抵达火星。它该做的事很简单：减速，进入环绕火星的轨道。地面控制人员发出指令，发动机点火，然后——信号消失了，再也没有回来。

事后调查委员会查遍了发动机、燃料、通信和软件，最后找到的原因让所有人哭笑不得：飞船没有坏，坏的是一次单位换算。制造飞船的团队在计算发动机推力时使用的是英制单位“磅力”，而导航团队的软件假定收到的数据是公制单位“牛”。这两个单位差着约 4.45 倍。于是每一次轨道修正，飞船实际被推的力道都比导航软件以为的大四倍多或小四倍多，误差一毫米一毫米地积累，九个多月后，飞船以为自己还在火星上空两百多公里的安全高度，实际上已经一头扎进了火星大气层，烧毁了。

这艘飞船的造价是三亿多美元。没有零件故障，没有程序死机，杀死它的是一个数字后面跟错了小尾巴。

如果你觉得航天事故离生活太远，再看一架民航客机。1983 年 7 月，加拿大航空 143 号航班——一架波音 767——在一万多米高空飞着飞着，两台发动机先后熄火。原因同样是单位：那架飞机是加拿大第一批用千克（而不是磅）计量燃料的机型，地勤按老习惯把“需要两万多千克”当成了“需要两万多磅”，实际只加了不到一半的油。万幸的是，机长恰好是一名滑翔机爱好者，硬是把这架没了动力的宽体客机滑翔了十几分钟，迫降在一条废弃跑道上，机上无人死亡。这架飞机后来得了一个绰号：“基米尼滑翔机”。同一个语法错误，在火星上烧掉三亿美元，在天上把三百多人推到鬼门关前。

你可能会说：那是工程师的事，跟我有什么关系？那我们换个近一点的场景。医生开的药片上写着“每片 0.25”，你敢不问单位就吞下去吗？0.25 克和 0.25 毫克差一千倍。再比如，把手机导航里“前方 500 米右转”改成“前方 500 右转”，这句话立刻变成废话。物理学里每一个有意义的陈述，都不是一个数，而是“一个数乘上一个公认的标准”。这一章开头请你记住这个悬念：为什么全世界最聪明的一群工程师，会栽在一件小学生都会的事情上？到本章结尾，我们会发现他们栽的不是算术，而是语法。

\step{先猜一猜}

在往下读之前，请先诚实地写下你对下面几个问题的直觉答案，不许翻书：

\begin{itemize}[leftmargin=1.8em]
  \item “一米”到底是谁规定的？如果明天国际计量大会投票把一米改长一倍，世界会怎样？
  \item 假如昨天夜里，宇宙中所有东西——包括你、你的米尺、地球和原子——都悄悄地膨胀了一倍，今天早上我们能发现吗？
  \item 一千克和一秒，哪一个更“基本”？能不能用秒来定义米？
  \item “光年”听起来像时间单位，它是吗？“千瓦·时”听起来像“每小时多少千瓦”，它是吗？
\end{itemize}

把你的猜测写下来。这一章结束时，我们会逐一回头看这些猜测。如果你现在觉得某些问题莫名其妙、根本无从猜起，那正好——说明你的直觉还没有为“单位”这件事建立过任何模型，而这正是本章要补上的东西。

\step{动手实验}

单位不是背出来的，是比较出来的。下面这个实验只需要一根细线、一把钥匙（或一个螺母）、一把卷尺和一部有秒表功能的手机。

\begin{experiment}[一把钥匙测出时间的语法]
第一步，感受“私人单位”的麻烦。先不用卷尺，用你的一扎（张开大拇指到中指的距离）量一量书桌的长度，记下来。再请家人用他们的一扎量同一张桌子。两个人的“桌长”几乎肯定不一样。你们都没有量错，错的是“扎”这个单位是私人的，拿出门就没法交流。再摸摸自己的脉搏，数一分钟跳多少下。你的“一分钟”是 75 跳，运动员可能只有 50 跳。看，连时间都有私人单位问题。

第二步，做一个摆。把钥匙系在细线一端，捏住线的另一端，让钥匙能自由摆动，这就做成一个单摆。拉开一个小角度放手，用手机秒表测它来回摆动 10 次的总时间，除以 10，得到一个来回的时间，叫周期。

第三步，改变摆长。分别取线长（从手指到钥匙中心）约 25 厘米、50 厘米、100 厘米，各测一次周期。你会发现：摆长变成 4 倍，周期并没有变成 4 倍，而是大约变成 2 倍；摆长变成 2 倍，周期大约变成 1.4 倍。

第四步，猜关系。什么样的数学关系能让“4 倍变成 2 倍”？开平方根：$4$ 的平方根是 $2$。于是可以大胆猜：周期的平方和摆长成正比，或者说周期正比于摆长的平方根。这个关系我们只靠线和钥匙就摸到了，但它还不够完整——周期还和什么有关？在月球上（重力只有地球的六分之一）同一个摆会变快还是变慢？先记下你的直觉，我们到“数学翻译器”一节用量纲这把尺子把答案逼出来。

第五步，检验一个反直觉的预言。多数人直觉上认为：挂重一点的东西，摆会被“坠”得更快。把两把钥匙系在同一条线上（质量翻倍），保持摆长不变，再测一次周期。你会发现周期几乎不变。请把这个“几乎不变”记在脑子里——到“数学翻译器”一节，我们会看到量纲分析早就预言了它：质量这个嫌疑人根本进不了公式。再顺手试一个变量：把拉开的角度从很小改成很大（接近平着放），周期会明显变长。这说明我们的猜测只在“小角度”下成立——物理公式常常附带这样的使用条件，本章“模型的边界”一节还会专门讨论它。
\end{experiment}

这个实验做了两件事：它先让你亲身体会“没有公认单位就没法交流”，再给你留下一个悬而未决的公式形状。请记住“4 倍变 2 倍”这个感觉，它是本章后面量纲分析的试验田。

\step{旧模型遇到麻烦}

关于单位，多数人的脑子里其实装着三个旧模型，它们在日常生活中大部分时候够用，所以一直没有被推翻。现在让它们挨个碰钉子。

\textbf{旧模型一：“测量结果就是一个数。”} 这个模型在买苹果时很好用：摊主说“三斤”，你掏钱，没人追问“三斤是什么”。但它遇到的历史麻烦可不小。中国古代的“一尺”，商代大约相当于今天的 17 厘米，汉代约 23 厘米，清代约 32 厘米。如果一位清代工匠按图纸给一座汉代遗址配一件“长三尺”的构件，他做出的东西会比原件长出将近一半。秦始皇统一度量衡之所以被写进史书，不是因为他喜欢整齐，而是因为“同一个数在不同地方意味着不同的量”会让税收、工程和贸易全部乱套。今天的国际贸易更极端：一船原油从美国运往欧洲，合同上写“桶”，结算按“吨”，炼油厂计量用“立方米”，三个数字背后是三套标准，任何一个环节把单位张冠李戴，损失都以百万美元计。测量结果从来不是“一个数”，而是“一个数乘以一个约定的标准”，少了标准，数就是废纸。

\textbf{旧模型二：“单位换算只是小学算术。”} 会背“1 米等于 100 厘米”不等于懂单位。火星气候轨道器的团队个个会背换算表，照样赔掉三亿美元。再看两个家里的例子。第一个是会让直觉出错的反例：你买了一块标称 500 GB 的移动硬盘，插到电脑上，系统却显示只有“465 GB”，是厂商骗人吗？不是。硬盘厂商按国际单位制，1 GB 等于 $10^9$ 字节；而许多操作系统内部按二进制计数，它显示的“GB”其实是 GiB，等于 $2^{30}$ 字节，约 $10.7$ 亿字节。500 GB 等于 $5\times10^{11}$ 字节，除以 $2^{30}$，正好是约 465 GiB。两边都没算错，用的是两本不同的“词典”。第二个例子是电费单。你家这个月用电 120“度”，一度电的正式名字叫“千瓦·时”。注意它是千瓦\textbf{乘}小时，不是千瓦\textbf{每}小时：一台 1000 瓦的电暖器开 1 小时，消耗的能量就是 1 千瓦乘 1 小时。它是个能量单位，换算成焦耳是 $3.6\times10^6$ 焦耳。把它念成“每小时多少千瓦”，就像把“平方米”念成“每米”一样，是语法错误。

\textbf{旧模型三：“单位随便选，反正能换算。”} 原则上对，实际上会付出惨重代价。真空光速约为 $3\times10^8$ 米每秒，换算成千米每小时是约 $1.08\times10^9$ 千米每小时——数字完全正确，也完全没有用，因为人对“十亿千米每小时”没有任何感觉。更深的问题是，选错单位会把自然规律打扮得面目可憎。后面第 9 章会讲到能量守恒，如果一边用焦耳、一边用卡、还混着千瓦·时，同一条定律写出来就满是丑陋的换算系数，丑到你会怀疑定律本身。单位的选择是一种“语言设计”：好的单位让规律显形，坏的单位让规律隐身。

三个旧模型倒下之后，露出来的是一个真空：我们从来没有认真回答过——一个物理量到底是什么？

\step{发明新概念}

现在像计量学家一样，把这个空缺填上。概念不多，只有三层。

\textbf{第一层：测量就是比较。} 说“桌子长 1.2 米”，意思不是桌子身上刻着“1.2”，而是：拿一个全世界公认的标准长度（一米），往桌边一比，桌边是这个标准的 1.2 倍。任何物理量都长成这个样子：
\[
\text{物理量} = \text{一个数} \times \text{一个公认的标准（单位）}.
\]
数是比例，单位是被拿来比较的那个“样品”。这句话可以检验：如果你写的某个“测量结果”拆不成这两部分，它就不是一个物理量。

\textbf{第二层：基本量与导出量。} 世界上的物理量成千上万，但让人惊讶的是，力学里只需要三种基本“样品”就能搭出全部：长度、质量、时间。国际单位制（SI）一共规定了七个基本量：长度（米）、质量（千克）、时间（秒）、电流（安培）、温度（开尔文）、物质的量（摩尔）、发光强度（坎德拉）。其余所有量的单位都可以用基本单位乘出来，叫导出单位。速度是米每秒，加速度是米每二次方秒，力的单位“牛”拆开就是千克乘米每二次方秒。这像搭积木：七块基本积木，搭出整个物理学的词汇表。

\textbf{第三层：量纲，物理量的“词性”。} 一个物理量由哪些基本量、各乘几次方构成，这叫它的量纲，记作方括号。长度的量纲记作 $[L]$，时间是 $[T]$，质量是 $[M]$。速度的量纲是 $[L]/[T]$，力是 $[M][L]/[T]^2$。这里要引进本章最重要的比喻，并且按规矩先说它像什么、再说它不像什么：量纲像语法里的词性——名词、动词、形容词。一个句子里“我吃桌子”之所以是病句，是因为“吃”这个动词后面跟了不该跟的词类；一个公式里“速度等于力”之所以是错句，是因为等号两边的词性对不上。但单位换算只像换方言：同样一个名词，普通话说“米”，换一种方言说“厘米”“英寸”，词性没变，意思也没变，只是发音不同。这个比喻不像的地方在于：人类语言的语法是人群约定俗成的，诗人可以故意违反语法制造美感；而量纲不是约定，它反映的是物理量本身的构造，公式两边量纲对不上就是铁板钉钉的错误，没有诗意可言。

最后讲一段单位定义的进化史，它能让你看清人类是怎么一步步把“标准”从手里挪到天上的。最初的一米来自地球：18 世纪末法国人规定，一米等于地球子午线（从北极到赤道那段）的一千万分之一。理想很宏伟，可地球表面在起伏、在变化，测绘也有误差，于是人们退而求其次，铸了一根铂铱合金的米原器，刻上两道线，说“这就是一米”。秒的故事类似：最初一秒是平太阳日的 $1/86400$，可地球自转正在受潮汐摩擦拖慢，今天的“一天”比恐龙时代长。拿一个正在变慢的东西当钟表的标准，等于用一把正在伸长的尺子量布。出路是一样的：到微观世界找不会变的过程。1967 年起，一秒被定义为铯原子一种特定振荡的 9192631770 个周期；1983 年起，一米被定义为光在真空中 $1/299792458$ 秒内走过的距离。

2019 年，这场迁徙终于完成：人类把七个基本单位全部改用自然常数来定义。在此之前，一千克的标准是巴黎郊外一个保险柜里的铂铱合金圆柱体，全世界的天平都要跟它比。可是金属会沾上灰尘、会磨损，它的质量万一变了，是全宇宙的千克跟着变，还是它自己“不合格”？这个逻辑漏洞困扰了人们一百多年。如今的定义反过来了：规定普朗克常数精确等于 $6.62607015\times10^{-34}$ 焦耳秒，规定真空光速精确等于 $299792458$ 米每秒，再由这些永不变的常数“生成”千克和米。标准不再是一件会生锈的文物，而是一条写进宇宙纹理里的定律。这件事本身就是一个深刻的物理宣言：我们相信最可靠的不是人造的样品，而是自然规律本身。为什么常数不会变、凭什么全宇宙的原子都遵守同一本账，这是后面量子篇章才能展开的问题，本章先把这个信念当作礼物收下。

\step{一句话模型}

\begin{oneline}
一个物理量 $=$ 一个数 $\times$ 一个公认的标准；单位不同只是换方言，量纲不同就是语病——所以看到任何公式，先检查等号两边的量纲能否对上。
\end{oneline}

\step{数学翻译器}

语法不能只用来欣赏，要能用来自动抓错。量纲分析就是这台抓错机器，它有三招，全部只需要初中数学。

\textbf{第一招：等号两边对账。} 等号是一条铁律：左边的量纲必须等于右边的量纲。以初中学过的匀加速运动位移公式为例：
\[
s = \frac{1}{2} a t^2 .
\]
左边 $s$ 是长度，量纲 $[L]$。右边 $a$ 的量纲是 $[L]/[T]^2$，$t^2$ 的量纲是 $[T]^2$，乘起来：$[L]/[T]^2 \times [T]^2 = [L]$。对上了。注意 $\frac{1}{2}$ 这种纯数字没有量纲，不参与对账。再故意写错一个：如果有人记成 $s = \frac{1}{2} a t$，右边量纲是 $[L]/[T]$，是个速度，跟左边的长度对不上——不用算任何数，这个公式已经被枪毙。同理，加减法也有规矩：只有量纲相同的量才能相加。“三米加五秒”不是等于八，而是根本无意义。

顺带练一手换算的基本功，它本质上也是对账。汽车时速 90 千米换成米每秒是多少？把单位当成分数里的因数来约：90 千米每小时 $= 90 \times \dfrac{1000\ \text{米}}{3600\ \text{秒}} = 25$ 米每秒。千米和米约掉，小时和秒约掉，剩下的就是答案。这类“链式约分”永远不会错，因为你每一步都在做量纲允许的运算。再看一个以后天天要见的量纲：能量。第 9 章会给出动能 $\frac{1}{2}mv^2$，它的量纲是 $[M][L]^2/[T]^2$；第 40 章会给出 $E = mc^2$，$c$ 是速度，所以右边量纲同样是 $[M][L]^2/[T]^2$。两个相隔三百年的公式说着同一种语法——这不是巧合，而是“能量”这个词在物理学里始终指着同一个东西。焦耳、卡、千瓦·时、电子伏特，全是这个词的不同方言。

\textbf{第二招：用特殊情形拷问公式。} 量纲对上了还不算完，马上拿零、极大、方向反转去试。还是 $s = \frac{1}{2} a t^2$：令 $t = 0$，得 $s = 0$——还没开始走，位移当然是零，合理；令 $a$ 翻倍，$s$ 跟着翻倍——推得越狠走得越远，合理；把 $a$ 取负号（刹车），$s$ 的增长变慢甚至变负——方向反转的物理意义也讲得通。一个公式如果经不起这种拷问，比如 $t = 0$ 时居然给出 $s = 5$ 米，那它多半还藏着没写出来的项。

\textbf{第三招：反过来，用量纲猜公式。} 这才是最神奇的一招。回到动手实验里的单摆。周期 $T$ 可能跟什么有关？列嫌疑人：摆球质量 $m$、摆长 $L$、还有决定“往下拉得多狠”的重力加速度 $g$（量纲 $[L]/[T]^2$）。摆动的角度先不管（实验时我们用了小角度）。现在问：$m$、$L$、$g$ 这三个量怎么组合，才能凑出一个量纲为 $[T]$ 的东西？质量 $m$ 的量纲里根本没有 $[T]$，它怎么乘都造不出时间来，所以 $T$ 跟 $m$ 无关——你没做实验的直觉可能正相反，这就是量纲分析的威力。剩下 $L$ 和 $g$：$L$ 除以 $g$ 的量纲是 $[L] \div ([L]/[T]^2) = [T]^2$，开个平方根就得到 $[T]$。所以答案被唯一地逼出来了：
\[
T = C\sqrt{\frac{L}{g}},
\]
其中 $C$ 是个无量纲的数，量纲分析算不出它。完整的理论（第 16 章会推导）给出 $C = 2\pi$。请立刻做特殊情形检查：$L$ 趋于零，$T$ 趋于零——摆没了，自然谈不上摆动周期；$g$ 越大，$T$ 越小——重力拉得越狠摆得越快，合理；到月球上 $g$ 只有地球六分之一，周期变成约 2.4 倍——回答了你实验时留下的问题。再对照你的实验数据：$L$ 从 25 厘米变到 100 厘米是 4 倍，$\sqrt{4} = 2$，周期正好该翻倍。一根线一把钥匙猜出的关系，和量纲分析逼出的公式，对上了。

\begin{mathbridge}
把上面“猜公式”的过程写成符号，可以看到它为什么唯一。设 $T \propto m^{\alpha} L^{\beta} g^{\gamma}$，两边取量纲：
\[
[T] = [M]^{\alpha}\, [L]^{\beta}\, \left([L][T]^{-2}\right)^{\gamma}
    = [M]^{\alpha}\, [L]^{\beta+\gamma}\, [T]^{-2\gamma}.
\]
比较两边各基本量的幂次，得三个方程：$\alpha = 0$，$\beta + \gamma = 0$，$-2\gamma = 1$。解得 $\alpha=0$，$\gamma=-\tfrac12$，$\beta=\tfrac12$，即 $T \propto \sqrt{L/g}$，别无他解。这套“列量纲方程组”的方法叫量纲分析，它在流体力学、天体物理里被用来在纸面上先“试算”出复杂现象的规律形状，再做实验验证系数。它的深层根据是：物理规律不应该依赖人类选用米还是英尺，因此公式两边必须在任何单位制下同时成立——这已经隐约碰到了第 57 章要讲的“对称性决定规律形式”的思想。
\end{mathbridge}

\textbf{数量级与科学记数法。} 量纲管“对不对”，数量级管“大不大”。物理学里的数跨度极大，统一用 10 的幂来写：光速 $3\times10^8$ 米每秒，原子直径约 $10^{-10}$ 米。所谓数量级，就是看 10 的头上的指数：$3\times10^8$ 和 $8\times10^8$ 算同一个数量级，$3\times10^8$ 和 $3\times10^5$ 差着三个数量级，也就是一千倍。估算时我们只关心指数，这正是费米估算的精髓。

\begin{center}
\begin{tikzpicture}[font=\small]
  \draw[-{Stealth}] (0,0) -- (12.8,0) node[right]{长度（米）};
  \foreach \x/\p/\name/\pos in {
    0/{10^{-10}}/{原子直径}/above,
    1.67/{10^{-5}}/{细胞}/below,
    3.33/{10^{0}}/{人的身高}/above,
    5.67/{10^{7}}/{地球直径}/below,
    7/{10^{11}}/{日地距离}/above,
    8.67/{10^{16}}/{一光年}/below,
    10.33/{10^{21}}/{银河系直径}/above,
    12/{10^{26}}/{可观测宇宙}/below}
    { \draw (\x,0.1) -- (\x,-0.1);
      \node[\pos=3pt, align=center] at (\x,0) {\name\\ $\p$}; }
\end{tikzpicture}
\end{center}

这张图按 10 的幂次等比排列，从原子到可观测宇宙横跨 36 个数量级。请你找到自己的位置，再找到头发丝（约 $5\times10^{-5}$ 米）和一粒米（约 $5\times10^{-3}$ 米）的位置。物理学的每个分支，其实都守在这把尺子的某一段上。

\textbf{带真实数据的估算：你一辈子心跳多少次？} 不看任何资料，只用量级常识：安静时心率约每分钟 75 次，一年有 $365\times24\times60 \approx 5\times10^5$ 分钟，寿命按 80 年即约 $4\times10^7$ 分钟。乘起来：$75 \times 4\times10^7 = 3\times10^9$ 次。三十亿次。这个数和医学观察的平均值在同一个数量级。再来一个费米式的问题：一座千万人口的大城市一天用多少自来水？每人每天饮用、做饭、洗澡、冲厕合计按 200 升估，$10^7$ 人乘 200 升是 $2\times10^9$ 升，即 200 万立方米，大约相当于一个边长 126 米的立方体水块。真实大城市日供水量确实在百万立方米量级。这种“先猜量级、再看结构”的功夫，比背任何公式都更接近物理学家的日常。

再练一题经典的：地球上的沙子多，还是天上的星星多？先估沙子：一粒沙直径按半毫米算，体积约 $10^{-10}$ 立方米；地球的海岸线、沙漠和浅海加起来，假设沙层总体积是 $10^{15}$ 立方米（这步最没把握，误差可能有两三个数量级，没关系）。沙子总数约 $10^{25}$ 粒。再估星星：银河系约有 $10^{11}$ 颗恒星，可观测宇宙里约有 $10^{11}$ 个星系，乘起来是 $10^{22}$ 颗。结论出人意料地稳：哪怕沙子的估计偏了一千倍，地球上每粒沙子对应的天上星星，也远不到一颗。费米估算的价值从不在答案精确，而在于告诉你哪些因素说了算、哪些因素随便猜也翻不了天。工程上也一样：芯片厂说“7 纳米制程”，就是在告诉你晶体管的尺寸守在 $10^{-8}$ 米那一段尺子上，离原子只剩两个数量级——这也是芯片越做越难的根本原因。

\begin{challenge}
在粒子物理里，理论家干脆把单位“烧掉”：规定真空光速 $c = 1$、约化普朗克常数 $\hbar = 1$。这意味着长度和时间用同一把尺量（$c$ 是长度除以时间，令它等于 1，时间就有了长度的量纲），质量、能量、动量也合并成同一种东西。于是质能关系从 $E = mc^2$ 简写成 $E = m$。这不是偷懒，而是一个深刻的陈述：秒和米、千克和焦耳之间的差别，本来就是人类度量习惯的差别，不是自然的差别。到第 38 章你会看到，时间和空间之所以能“混用同一把尺”，根源在相对论；到第 40 章，$E=mc^2$ 本身会被重新推导。本章只要记住：单位是人为的，量纲背后的结构是自然的，而某些“换算系数”（如 $c$）其实是自然常数，它们的存在暗示着被我们当成两回事的量，可能本是同一件事。
\end{challenge}

\step{模型的边界}

量纲分析是抓错机器，但它有明确的盲区，用过头就会翻车。请牢记四条边界。

\textbf{边界一：量纲对，不等于公式对。} 单摆公式里的 $2\pi$、位移公式里的 $\frac{1}{2}$，量纲分析一概看不见。更隐蔽的是无量纲的函数：假如单摆周期实际上还乘着一个随摆角变化的因子（事实上的确如此，第 16 章会讲），量纲分析也无能为力，因为摆角是无量纲量。正确说法是：量纲错误必然公式错误；量纲正确只是拿到了准考证，不是录取通知书。

\textbf{边界二：别拿量纲分析代替物理。} 有这样一个真实教训式的例子：雨滴从 1000 米高空落下，用 $v = \sqrt{2gh}$ 可以算出落地速度约 140 米每秒，单位检查完全通过。可真实雨滴落地只有几米每秒——公式语法没错，错在物理模型漏掉了空气阻力。量纲分析检查的是“这句话通不通”，不检查“这句话真不真”。第 1 章说过，模型对不对，最终要实验说了算。

\textbf{边界三：有些量天生没有量纲。} 角度（用弧度时）、折射率、百分比、概率，都是数与数的比值，单位约掉了。它们可以和任何量纲的量相乘而不留痕迹，这正是无量纲系数能躲过量纲检查的原因，也是它们容易被人乱用的原因——“效率 120\%”在语法上毫无毛病，在热力学里却是第 30 章要宣判的死刑。

还有一类麻烦来自“名字像、量纲不同”的邻居。充电宝上印着 10000 mAh（毫安时）：毫安是电流，时是时间，电流乘时间是电荷量，不是能量。要知道它存了多少能量，还得乘上电池的电压——电压的物理意义正是“每份电荷携带多少能量”，这一点第 20 章会仔细讲。商家有意无意地利用这种混淆做广告，而量纲分析能让你一眼看穿：只报电荷不报电压，就像只报“一桶水”不报从多高的地方倒下来。

\textbf{边界四：单位系统本身也有失效的尺度。} 把长度一路除以 10，除到约 $1.6\times10^{-35}$ 米（普朗克长度），把时间除到约 $5.4\times10^{-44}$ 秒（普朗克时间），目前的物理理论（广义相对论与量子力学）会同时失效，那里“长度”“时间”这两个基本量本身还有没有意义，没有人知道。这是第 41 章的伏笔。也就是说，本章这套语法管用的范围虽然横跨 36 个数量级，但它仍然有边疆——承认边疆的存在，恰恰是科学的诚实。

\step{回头解释开场现象}

现在，用量纲这副眼镜重新看火星气候轨道器事故。

发动机做轨道修正时，真正决定飞船轨迹改变量的物理量是冲量：推力乘以作用时间。导航软件每秒都在收一个数——发动机报了多大冲量——并把它当成“牛·秒”代入轨道计算。而发动机的报表写的是“磅力·秒”。1 磅力约等于 4.45 牛，所以软件每一次都把真实的冲量读错了 4.45 倍。用本章的话说：两边交换的是没有单位的裸数字，量纲对账这道工序被省略了。等式的一边是公制的牛顿秒，另一边是英制的磅力秒，词性（量纲）虽然碰巧相同，方言却不同，而软件连方言转换都没做。误差不大——每次只有几倍——但飞船在九个多月里反复修正，每次都带着同一个系统性的错误，轨迹被一点一点地“拽”向火星，最终在进入高度上差了约一百多公里，撞进了大气层。

事后的整改建议里，最朴素的一条恰恰最像本章的习题：所有接口数据必须显式携带单位，软件在代入计算前先做单位一致性检查。三亿多美元买到的教训，就是本章开头那句话——一个物理量等于一个数乘以一个公认的标准，裸奔的数字不配上飞船。

顺便也回答“先猜一猜”里的问题。一米是谁定的？曾经是一根金属棒，如今是“光在真空中 $1/299792458$ 秒内走的距离”——用秒和自然常数定义米，所以秒比米更“基本”。如果万物（包括米尺和原子）同步膨胀一倍，我们原则上无从发现：测量是比较，所有比较结果都不变——这说明物理学里真正有意义的是量与量之间的比值，裸尺度本身反而可能是多余的描述，这个念头会在第 38 章变得非常重要。光年不是时间，是光跑一年的距离，约 $9.5\times10^{15}$ 米；千瓦·时不是“每小时千瓦”，是功率乘时间，是能量。你猜对了几个？

\step{脑洞实验与挑战题}

\begin{thought}[怎么向外星人解释“一米”]
1972 年发射的先驱者 10 号探测器携带了一块镀金铝板，要送给可能截获它的外星文明。设计者们立刻撞上本章的核心困难：你不能在板上刻“此人身高 1.7 米”，因为外星人的“米”和我们的米毫无关系，一切人为单位都无法翻译。他们的解决方案漂亮极了：不在板上写任何人为单位，只画图。板上刻着氢原子自旋翻转时发出的辐射——这是全宇宙处处相同的物理过程——它的波长约为 21 厘米，周期约为 $7\times10^{-10}$ 秒。外星物理学家只要能认出这个图案，就同时拿到了一把尺子和一只钟：板上人体的身高按这把尺子是 8 个波长，一目了然。这个思想实验说明：单位可以是人造的，但物理量之间的比例是宇宙通用的语言。至于氢原子为什么会发出波长如此确定的辐射、为什么全宇宙的氢都一样，那是第 51、52 章要揭开的量子世界的故事。现在再追一步：假如未来人类发现某个“自然常数”其实在缓慢变化，这套以常数为基础的单位制会发生什么？这不是科幻，而是当代宇宙学真的在检验的问题。
\end{thought}

\begin{puzzles}
\item 一位同学推导自由落体，得到结论“下落距离 $h = v^2 g / 2$”（$v$ 是末速度，$g$ 是重力加速度）。不替他重算，你能立刻断定他错了吗？正确的关系应该长什么样？（提示：分别写出等号两边的量纲，再比较；然后把正确的量纲组合凑出来，看看系数 $\frac{1}{2}$ 你能不能凭量纲确定。）
\item 宇航员想比较月球表面和地球表面的重力加速度谁大谁小，手里只有一根线、一个螺母、一把卷尺和一只秒表。他可以怎样做？需要知道线长的精确数值吗？（提示：本章的单摆公式里，$L$ 和 $g$ 以什么方式捆绑出现？两次测量取什么比值可以消掉误差？）
\item 估算：把全校学生的头发（假设每人留着 10 厘米长的头发）一根接一根连起来，能绕地球赤道多少圈？请写出你每一步的假设和数量级，最后只报 10 的幂。（提示：地球周长可以从本章的尺度图里找；人数自己大胆估；重点不是答案，是你敢不敢在信息不全时往下走。）
\item 假如某国立法规定：从此时间改用“心跳”（平均一次心跳为一单位），长度改用“步幅”。物理定律会被这项法律改变吗？哪些东西会变，哪些东西不变？（提示：区分“规律的数值形式”和“规律本身”；想想单摆周期公式里的 $g$，它的数值会怎样，它的量纲又会怎样。）
\end{puzzles}

这一章没有教你任何一条新定律，却给了你一件能用一辈子的工具：看任何公式，先问它每个符号的量纲；听任何数字，先问它后面跟着什么单位。从下一章开始，我们要让物理量动起来——位置怎样随时间变化，图像会替我们说话。届时你会发现，坐标轴上写的每一个标签，都必须遵守今天立下的语法。
